% GENERATED FILE. Edit canonical lessons, then run npm run generate:books.
% canonical-source-hash: fnv1a64:e68d2da9a1bcf5cd
% canonical-lessons: GU-C294-lapasvu, GU-C294-ujhardo, GU-C294-gabadvu, GU-C294-dhal, GU-C294-bhan

\chapter{A Fall in the Rain}
\label{ch:gu-a-fall-in-the-rain}
\hlchaptermodality{294}

\begin{chapteropening}
\textbf{By the end of this chapter:} \emph{I can say how someone slipped and fell and when they came round.}

Complete the last lesson of chapter 294: I can say how someone slipped and fell and when they came round.
\end{chapteropening}

\section[\texorpdfstring{\gu{લપસવું} (lapasvũ)}{લપસવું (lapasvũ)}]{\gu{લપસવું} (lapasvũ) — to slip}

\label{lesson:GU-C294-lapasvu}



\begin{quote}
Before the new one: say the Gujarati for serious, then the Gujarati for to panic.
\end{quote}

\subsection*{You'll want to know: \gu{લપસવું}}
\textbf{\gu{લપસવું}} — \emph{lapasvũ} — \textquotedblleft{}to slip\textquotedblright{}.

\begin{quote}
\gu{વરસાદમાં} \gu{મીરા} \gu{પગથિયાં} \gu{પર} \gu{લપસી} \gu{ગઈ}.
\end{quote}

\begin{quote}
\emph{varsādmā̃ Mīrā pagthiyā̃ par lapsī gaī.}
\end{quote}

\begin{quote}
In the rain Mira slipped on the steps.
\end{quote}

Past (she …): \gu{તે} \gu{લપસી} \gu{ગઈ}, \emph{te lapsī gaī}.

\begin{cousinweb}
\textbf{\gu{લપસવું}} (\emph{lapasvũ}) No dictionary we use records an older form, so its history is not settled. A wet floor makes you \textbf{\gu{લપસી} \gu{જવું}}: \textbf{\gu{ભીની} \gu{ભોંય} \gu{પર} \gu{લપસી} \gu{ગઈ}}, she slipped on the wet floor.
\end{cousinweb}

\begin{grammarlens}[title={What you've built}]
One more word for telling what happened.
\end{grammarlens}

\subsection*{Your turn}
\begin{itemize}
\raggedright
  \item \emph{Say it:} \emph{lapasvũ}
  \item \emph{Say it:} \emph{lapasvũ}, once more
  \item \emph{Say it:} \emph{varsādmā̃ Mīrā pagthiyā̃ par lapsī gaī.} — In the rain Mira slipped on the steps.
  \item \emph{Recall:} say the Gujarati for serious, then the Gujarati for to panic, then say \emph{lapasvũ} again
\end{itemize}

\begin{tcolorbox}[breakable,colback=teal!4,colframe=teal!35!black,title={Before you move on}]
What does \gu{લપસવું} mean? (\textquotedblleft{}To slip\textquotedblright{}.) Say it once more.
\end{tcolorbox}
\section[\texorpdfstring{\gu{ઉઝરડો} (ujharḍo)}{ઉઝરડો (ujharḍo)}]{\gu{ઉઝરડો} (ujharḍo) — a graze; a scratch}

\label{lesson:GU-C294-ujhardo}



\begin{quote}
Before the new one: say the Gujarati for to panic, then the Gujarati for to slip.
\end{quote}

\subsection*{You'll want to know: \gu{ઉઝરડો}}
\textbf{\gu{ઉઝરડો}} — \emph{ujharḍo} — \textquotedblleft{}a graze; a scratch\textquotedblright{}.

\begin{quote}
\gu{મીરાના} \gu{ઘૂંટણ} \gu{પર} \gu{મોટો} \gu{ઉઝરડો} \gu{હતો}.
\end{quote}

\begin{quote}
\emph{Mīrānā ghū̃ṭaṇ par moṭo ujharḍo hato.}
\end{quote}

\begin{quote}
Mira had a big graze on her knee.
\end{quote}

A masculine noun.

\begin{cousinweb}
\textbf{\gu{ઉઝરડો}} (\emph{ujharḍo}) comes from \emph{*ujjharati}, a form CDIAL reconstructs, and CDIAL lists Gujarati \emph{ujharṛo}, \textquotedblleft{}a scratch\textquotedblright{} (CDIAL 1675). It is skin scraped in a fall: \textbf{\gu{ઘૂંટણ} \gu{પર} \gu{ઉઝરડો}}.
\end{cousinweb}

\begin{grammarlens}[title={What you've built}]
One more word for telling what happened.
\end{grammarlens}

\subsection*{Your turn}
\begin{itemize}
\raggedright
  \item \emph{Say it:} \emph{ujharḍo}
  \item \emph{Say it:} \emph{ujharḍo}, once more
  \item \emph{Say it:} \emph{Mīrānā ghū̃ṭaṇ par moṭo ujharḍo hato.} — Mira had a big graze on her knee.
  \item \emph{Recall:} say the Gujarati for to panic, then the Gujarati for to slip, then say \emph{ujharḍo} again
\end{itemize}

\begin{tcolorbox}[breakable,colback=teal!4,colframe=teal!35!black,title={Before you move on}]
What does \gu{ઉઝરડો} mean? (\textquotedblleft{}A graze; a scratch\textquotedblright{}.) Say it once more.
\end{tcolorbox}
\section[\texorpdfstring{\gu{ગબડવું} (gabaḍvũ)}{ગબડવું (gabaḍvũ)}]{\gu{ગબડવું} (gabaḍvũ) — to tumble}

\label{lesson:GU-C294-gabadvu}



\begin{quote}
Before the new one: say the Gujarati for to slip, then the Gujarati for a graze.
\end{quote}

\subsection*{You'll want to know: \gu{ગબડવું}}
\textbf{\gu{ગબડવું}} — \emph{gabaḍvũ} — \textquotedblleft{}to tumble\textquotedblright{}.

\begin{quote}
\gu{છોકરો} \gu{સાઈકલ} \gu{પરથી} \gu{ગબડી} \gu{પડ્યો}.
\end{quote}

\begin{quote}
\emph{chokro sāīkal parthī gabḍī paḍyo.}
\end{quote}

\begin{quote}
The boy tumbled off his bicycle.
\end{quote}

Past (he …): \gu{તે} \gu{છોકરો} \gu{ગબડી} \gu{પડ્યો}, \emph{te chokro gabḍī paḍyo}.

\begin{cousinweb}
\textbf{\gu{ગબડવું}} (\emph{gabaḍvũ}) No dictionary we use records an older form, so its history is not settled. It is falling and rolling over: \textbf{\gu{ગબડી} \gu{પડ્યો}}, he went tumbling down.
\end{cousinweb}

\begin{grammarlens}[title={What you've built}]
One more word for telling what happened.
\end{grammarlens}

\subsection*{Your turn}
\begin{itemize}
\raggedright
  \item \emph{Say it:} \emph{gabaḍvũ}
  \item \emph{Say it:} \emph{gabaḍvũ}, once more
  \item \emph{Say it:} \emph{chokro sāīkal parthī gabḍī paḍyo.} — The boy tumbled off his bicycle.
  \item \emph{Recall:} say the Gujarati for to slip, then the Gujarati for a graze, then say \emph{gabaḍvũ} again
\end{itemize}

\begin{tcolorbox}[breakable,colback=teal!4,colframe=teal!35!black,title={Before you move on}]
What does \gu{ગબડવું} mean? (\textquotedblleft{}To tumble\textquotedblright{}.) Say it once more.
\end{tcolorbox}
\section[\texorpdfstring{\gu{ઢાળ} (ḍhāḷ)}{ઢાળ (ḍhāḷ)}]{\gu{ઢાળ} (ḍhāḷ) — a slope}

\label{lesson:GU-C294-dhal}



\begin{quote}
Before the new one: say the Gujarati for a graze, then the Gujarati for to tumble.
\end{quote}

\subsection*{You'll want to know: \gu{ઢાળ}}
\textbf{\gu{ઢાળ}} — \emph{ḍhāḷ} — \textquotedblleft{}a slope\textquotedblright{}.

\begin{quote}
\gu{આ} \gu{રસ્તા} \gu{પર} \gu{મોટો} \gu{ઢાળ} \gu{છે}.
\end{quote}

\begin{quote}
\emph{ā rastā par moṭo ḍhāḷ che.}
\end{quote}

\begin{quote}
This road has a steep slope.
\end{quote}

A masculine noun.

\begin{cousinweb}
\textbf{\gu{ઢાળ}} (\emph{ḍhāḷ}) probably belongs with \emph{ḍhaḷvũ}, to slip, to be poured out, which CDIAL sets under \emph{*ḍhalati} (CDIAL 5581): ground that falls away. A bicycle goes fast \textbf{\gu{ઢાળ} \gu{પર}}, on a slope.
\end{cousinweb}

\begin{grammarlens}[title={What you've built}]
One more word for telling what happened.
\end{grammarlens}

\subsection*{Your turn}
\begin{itemize}
\raggedright
  \item \emph{Say it:} \emph{ḍhāḷ}
  \item \emph{Say it:} \emph{ḍhāḷ}, once more
  \item \emph{Say it:} \emph{ā rastā par moṭo ḍhāḷ che.} — This road has a steep slope.
  \item \emph{Recall:} say the Gujarati for a graze, then the Gujarati for to tumble, then say \emph{ḍhāḷ} again
\end{itemize}

\begin{tcolorbox}[breakable,colback=teal!4,colframe=teal!35!black,title={Before you move on}]
What does \gu{ઢાળ} mean? (\textquotedblleft{}A slope\textquotedblright{}.) Say it once more.
\end{tcolorbox}
\section[\texorpdfstring{\gu{ભાન} (bhān)}{ભાન (bhān)}]{\gu{ભાન} (bhān) — consciousness}

\label{lesson:GU-C294-bhan}



\begin{quote}
Before the new one: say the Gujarati for to tumble, then the Gujarati for a slope.
\end{quote}

\subsection*{You'll want to know: \gu{ભાન}}
\textbf{\gu{ભાન}} — \emph{bhān} — \textquotedblleft{}consciousness\textquotedblright{}.

\begin{quote}
\gu{અડધા} \gu{કલાક} \gu{પછી} \gu{તે} \gu{ભાનમાં} \gu{આવ્યો}.
\end{quote}

\begin{quote}
\emph{aḍdhā kalāk pachī te bhānmā̃ āvyo.}
\end{quote}

\begin{quote}
Half an hour later he came round.
\end{quote}

A neuter noun.

\begin{cousinweb}
\textbf{\gu{ભાન}} (\emph{bhān}) is Sanskrit \emph{bhāna}, \textquotedblleft{}appearance, perception\textquotedblright{} (Monier-Williams): in Gujarati, being aware of oneself and the world. \textbf{\gu{ભાનમાં} \gu{આવ્યો}}, he came round; \textbf{\gu{બેભાન}} is its opposite.
\end{cousinweb}

\begin{grammarlens}[title={What you've built}]
One more word for telling what happened.
\end{grammarlens}

\subsection*{Your turn}
\begin{itemize}
\raggedright
  \item \emph{Say it:} \emph{bhān}
  \item \emph{Say it:} \emph{bhān}, once more
  \item \emph{Say it:} \emph{aḍdhā kalāk pachī te bhānmā̃ āvyo.} — Half an hour later he came round.
  \item \emph{Recall:} say the Gujarati for to tumble, then the Gujarati for a slope, then say \emph{bhān} again
\end{itemize}

\begin{tcolorbox}[breakable,colback=teal!4,colframe=teal!35!black,title={Before you move on}]
What does \gu{ભાન} mean? (\textquotedblleft{}Consciousness\textquotedblright{}.) Say it once more.
\end{tcolorbox}
